\section{Generalized diffusion：不同公式，其实同一思想}

\subsection{统一写法}

到这里，diffusion 的不同流派已经足够多了。lecture 里进一步给了一个统一写法。把不同版本的 noising path 和预测目标都写成同一个模板后，很多看起来不同的算法会突然变得可比：

\[
x_t = a(t) x + b(t) z,\qquad
y_{\text{gt}} = c(t)x + d(t)z,\qquad
y_{\text{pred}} = f_\theta(x_t,t).
\]

训练目标就是让 $y_{\text{pred}}$ 逼近 $y_{\text{gt}}$。

\begin{figure}[H]
\centering
\includegraphics[width=0.95\textwidth]{slides-images/slide-108.jpg}
\caption{Generalized diffusion 的统一形式：rectified flow 只是其中一个特例\protect\footnotemark}
\end{figure}
\footnotetext{Slides 第109页。}

\begin{figure}[H]
\centering
\includegraphics[width=0.95\textwidth]{slides-images/slide-109.jpg}
\caption{把 rectified flow 写成通用模板后，就能看清它和其他 diffusion 变体的关系\protect\footnotemark}
\end{figure}
\footnotetext{Slides 第110页。}

\subsection{几个最重要的特例}

在这个统一模板里，不同方法只是对 $a,b,c,d$ 的选择不同：

\begin{figure}[H]
\centering
\includegraphics[width=0.95\textwidth]{slides-images/slide-110.jpg}
\caption{Rectified flow 作为特例：$a(t)=1-t,\ b(t)=t,\ c(t)=-1,\ d(t)=1$\protect\footnotemark}
\end{figure}
\footnotetext{Slides 第111页。}

\begin{figure}[H]
\centering
\includegraphics[width=0.95\textwidth]{slides-images/slide-111.jpg}
\caption{Variance preserving (VP) 的思路：保持中间态的方差结构更稳定\protect\footnotemark}
\end{figure}
\footnotetext{Slides 第112页。}

\begin{figure}[H]
\centering
\includegraphics[width=0.95\textwidth]{slides-images/slide-112.jpg}
\caption{Variance exploding (VE) 的思路：让噪声足够大，彻底淹没原始信号\protect\footnotemark}
\end{figure}
\footnotetext{Slides 第113页。}

\begin{figure}[H]
\centering
\includegraphics[width=0.95\textwidth]{slides-images/slide-113.jpg}
\caption{预测目标的几种写法：x-prediction、epsilon-prediction 和 v-prediction\protect\footnotemark}
\end{figure}
\footnotetext{Slides 第114页。}

\begin{knowledgebox}{为什么这些写法本质上是同一件事}
\begin{itemize}
    \item x-prediction 直接预测干净样本
    \item epsilon-prediction 预测噪声
    \item v-prediction 预测二者的线性组合
\end{itemize}
它们看起来不同，但都在学同一条“从噪声回到数据”的路径，只是参数化方式不同。
\end{knowledgebox}

\subsection{Diffusion 的三种解释}

lecture 进一步指出，diffusion 不只是“一个训练技巧”，它还有三种常见解释。

\begin{figure}[H]
\centering
\includegraphics[width=0.95\textwidth]{slides-images/slide-115.jpg}
\caption{Diffusion 作为 latent variable model：前向加噪是已知过程，反向生成是学出来的近似后验\protect\footnotemark}
\end{figure}
\footnotetext{Slides 第116页。}

第一种是 latent variable model 视角：前向过程把数据变成 latent trajectory，反向过程则通过 ELBO 或相关目标去拟合。第二种是 score matching 视角：模型学习的是 $\nabla_x \log p(x)$，也就是指向高密度区域的向量场。第三种是 SDE / ODE 视角：噪声演化可以写成连续时间动力系统，采样就是数值积分。

\begin{figure}[H]
\centering
\includegraphics[width=0.95\textwidth]{slides-images/slide-116.jpg}
\caption{Diffusion 学到的是 score：梯度场把样本往高概率区域推\protect\footnotemark}
\end{figure}
\footnotetext{Slides 第117页。}

\begin{figure}[H]
\centering
\includegraphics[width=0.95\textwidth]{slides-images/slide-117.jpg}
\caption{SDE 视角：连续 noising 可以写成微分方程，采样就是沿该动态系统反向积分\protect\footnotemark}
\end{figure}
\footnotetext{Slides 第118页。}

\begin{warningbox}{这三种视角不是互斥的}
把 diffusion 解释成 latent model、score model 或 SDE，并不意味着它们是三种不同算法。更准确的说法是：它们是同一套生成过程的不同投影。工程上，你只需要选一个最方便实现和调参的参数化。
\end{warningbox}

\begin{figure}[H]
\centering
\includegraphics[width=0.95\textwidth]{slides-images/slide-118.jpg}
\caption{一个很好的外部阅读线索：Sander Dieleman 关于 diffusion 视角的总结博客\protect\footnotemark}
\end{figure}
\footnotetext{Slides 第119页。}

